\chapter{Le modèle de menace}

Avant de lister des défenses, il faut dire ce qu'on défend et contre qui. Sans
cela on empile des mécanismes à la mode sans savoir lesquels comptent.

\section*{Ce que la plateforme expose}

InsaCloud n'est pas un site vitrine. Elle expose délibérément trois choses que
l'on ne met normalement pas sur un réseau :

\begin{itemize}[leftmargin=14pt,itemsep=3pt]
\item \textbf{Des comptes \cmd{root}}, un par machine louée, avec leur mot de
      passe, remis à des utilisateurs qu'on ne connaît pas.
\item \textbf{Une plage de ports ouverts} --- 8000 à 9000 --- derrière lesquels
      répondent des serveurs SSH, des terminaux web et des bureaux graphiques.
\item \textbf{Un démon Docker} sur chaque nœud, c'est-à-dire un service dont le
      contrôle équivaut à être \cmd{root} sur la machine hôte.
\end{itemize}

\begin{notion}{Pourquoi le démon Docker est l'actif le plus sensible}
Quiconque peut parler au démon Docker peut lancer un conteneur privilégié qui
monte le disque de l'hôte. Il n'y a pas d'escalade à trouver : c'est le
fonctionnement normal de Docker. « Accès au démon Docker » et « accès
administrateur de la machine » sont donc deux formulations de la même chose, et
toute décision d'architecture doit partir de là.
\end{notion}

\section*{Les quatre adversaires}

\begin{center}
\small
\begin{tabularx}{\textwidth}{@{}p{3.6cm}X@{}}
\toprule
\textbf{Qui} & \textbf{Ce qu'il cherche, et ce qui l'arrête} \\
\midrule
\textbf{Le robot d'Internet} &
Il balaie les ports, tente des mots de passe SSH et des dizaines de milliers de
chemins HTTP connus. C'est l'adversaire le plus \emph{fréquent} et le plus
bête. Parades : authentification par clé, limitation des tentatives, en-têtes
et politique de contenu stricts. \\
\addlinespace
\textbf{L'utilisateur curieux} &
Un étudiant légitime qui essaie de voir les machines d'un autre, de prolonger
sa location sans payer le temps, ou de sortir de son conteneur. C'est
l'adversaire le plus \emph{probable} dans un contexte pédagogique. Parades :
cloisonnement par propriétaire vérifié en base, capacités minimales dans le
conteneur, quotas. \\
\addlinespace
\textbf{Le locataire malveillant} &
Il loue une machine en règle, puis s'en sert comme tremplin : scanner le réseau
interne, atteindre le contrôleur, miner, attaquer l'extérieur. Il est
\emph{déjà \cmd{root}} dans son conteneur --- c'est le service rendu. Parades :
durcissement du conteneur, réseau dédié, limites de ressources, durée de vie
courte. \\
\addlinespace
\textbf{Celui qui lit le dépôt} &
Le code est sur GitHub. Un secret commis par erreur, même retiré ensuite, reste
dans l'historique. Parades : \cmd{gitleaks} sur tout l'historique à chaque
envoi, Ansible Vault, fichiers de clés exclus du dépôt. \\
\bottomrule
\end{tabularx}
\end{center}

\section*{Ce qui est hors périmètre, et pourquoi}

Un modèle de menace honnête nomme aussi ce qu'il abandonne.

\textbf{Un attaquant qui a déjà \cmd{root} sur le contrôleur} a gagné : il lit
la base, les clés, et pilote tous les nœuds. Aucune des mesures décrites ici ne
le gêne. La défense contre ce cas est organisationnelle (qui a accès à la
machine physique), pas logicielle.

\textbf{Une évasion de conteneur par une faille du noyau.} Les conteneurs
partagent le noyau de l'hôte ; une vulnérabilité du noyau les traverse. Le
durcissement réduit la surface (capacités retirées,
\cmd{no-new-privileges}) mais ne crée pas la frontière qu'offrirait une machine
virtuelle. C'est le compromis assumé de la conteneurisation, et l'énoncé le
tranche en imposant Docker.

\textbf{Le déni de service distribué.} Les limites de ressources par conteneur
protègent les \emph{voisins} d'un locataire gourmand ; elles ne protègent pas la
plateforme d'une attaque volumétrique venue de l'extérieur, qui relève de
l'hébergeur.

\begin{retenir}
Trois actifs à protéger : le \textbf{démon Docker} (équivalent à \cmd{root} sur
l'hôte), les \textbf{secrets} (clés, mots de passe des machines), et le
\textbf{cloisonnement entre locataires}. Tout le reste en découle.
\end{retenir}

\section*{Les dix couches de défense}

Le reste du document suit l'ordre d'une requête, du navigateur jusqu'au
conteneur :

\begin{center}
\small
\begin{tabularx}{\textwidth}{@{}p{0.6cm}p{4.2cm}X@{}}
\toprule
\textbf{\#} & \textbf{Couche} & \textbf{Ce qu'elle empêche} \\
\midrule
2 & Identité et authentification & qu'on se fasse passer pour un utilisateur \\
3 & Session et navigateur & qu'on détourne une session ouverte \\
4 & Secrets & qu'un secret traîne en clair quelque part \\
5 & Accès aux machines louées & qu'on entre dans la machine d'un autre \\
6 & Durcissement des conteneurs & qu'un locataire sorte de sa machine \\
7 & Contrôleur et nœuds & qu'on prenne la main sur un démon Docker \\
8 & Données et réseau & qu'on corrompe l'état ou qu'on écoute le trafic \\
9 & Chaîne d'intégration & qu'une faille entre dans le dépôt sans qu'on le voie \\
10 & Limites assumées & qu'on se croie protégé là où on ne l'est pas \\
\bottomrule
\end{tabularx}
\end{center}
